\honest{Lonely Dissent}{Eliezer Yudkowsky}{2007}

Asch's experiment, I say, showed that a single ally makes dissent easy, and that being the first to dissent is ``a lot harder.'' \nb{Both are true. Asch also found that even subjects facing the whole group alone gave the correct answer about two times in three, and that a quarter of them never gave in at all. The post leaves this out.}

Lonely dissent, I say, feels like going to school in a clown suit, not like going to school dressed in black. I say: ``If there's one thing I can't stand, it's fakeness.'' Lonely dissent, I add, is one of the most often faked traits, since everyone wants to be an iconoclast.

Vegetarians show some courage, I expect, but people understand them. Wearing black to school is a standard rebellion that everyone recognizes. Real courage means facing people who simply find you weird. My example is explaining cryonics, the freezing of the dead in the hope of reviving them later, which I support. \nb{The imagined listeners do doubt that it will work (``You think that's going to stop you from dying?''), but the post treats their doubt as incomprehension, not as a judgment that might be right.}

Then an evolutionary story, which I say I am ``tempted'' to offer as a post facto explanation. A small group could leave the tribe together, but a person driven out alone probably died. The case of cryonics, I say, shows that ``the fear of thinking really different is stronger than the fear of death,'' because ``many more of us go skydiving than sign up for cryonics.'' \nb{A skydive is one ticket with a known result. Cryonics costs tens of thousands of dollars, and no one has ever been revived. The post does not mention either difference.} Twice, in consecutive paragraphs, I explain that we do not reason this out explicitly, because ``that is not the nature of evolutionary psychology'' and because ``that's not how evolutionary psychology works.''

The highest courage, I say, belongs to the first person to say a new thing, such as Robert Ettinger, who founded cryonics. From him I move to scientific revolutionaries, with every qualification: revolution is not the only route to greatness, you cannot get there by trying, and you need all existing knowledge plus luck. \nb{The post links Ettinger and the scientific revolutionaries by one feature, being first. It does not ask whether cryonics is also right, which a few paragraphs later it calls the hard part.}

Then I say that not every dissenting idea is good, and that ``Most of the difficulty in having a new true scientific thought is in the `true' part.'' So differ only when you have an overwhelmingly good reason. In a footnote I cite Robin Hanson's post ``Against Free Thinkers.'' \nb{Hanson's post argues that ``undiscriminating freethinkers are our main obstacle to innovation.'' The post does not apply this to cryonics.}

Then I say: ``Now me, you know, I really am an iconoclast.'' Everyone thinks they are, I say, with irony, but in my case it is true: I would ``totally'' have worn the clown suit. Then, seriously, I say that dissenting too easily is my own nature, and that I have to correct for it. I close: ``You wouldn't want to end up as a free thinker.''
